\part{Telemetry, monitoring and feature engineering}

\section{Telemetry collection}

\begin{qa}{264}{What network metrics does CARS-SF observe at runtime?}
Per 1\,s tick (\code{telemetry/collector.py}):
\begin{itemize}
\item per-service offered uplink and delivered downlink (access-port counters);
\item uplink, downlink and max utilisation of the bottleneck;
\item backhaul and backup throughput;
\item HTB class sent bytes and packets, drops and backlog, per direction;
\item per-service RTT (ICMP probes);
\item link carrier state and timestamped transitions (50\,ms watcher);
\item OVS flow bytes and counts (OVS only);
\item host CPU and memory.
\end{itemize}
\end{qa}

\begin{qa}{265}{Which metrics come from OVS counters, tc, active probes or application-level measurement?}
OVS: flow bytes and counts. tc: per-class bytes, drops, backlog. Probes: RTT and reachability. Kernel interface
counters: rates and link state. Application-level (per-flow one-way delay, loss, throughput): only post-run from
captures, plus live captures for the in-loop classifier.
\end{qa}

\begin{qa}{266}{Why is a single telemetry source insufficient?}
Counters give rates but no delay. Probes give delay but no per-class volume. HTB stats show drops but not their
cause upstream. Link state shows cuts but not degradation. The predictor needs RTT trends \emph{and} local queue
evidence to tell congestion from backhaul degradation; that confusion caused a false failover on OVS until both
signals were combined.
\end{qa}

\begin{qa}{267}{How often is telemetry collected, and why that interval?}
\begin{itemize}
\item 1\,s ticks: enough to see 5\,s control decisions, cheap, and well above timer noise.
\item 50\,ms link polling: carrier loss must be acted on in under 1\,s.
\item Probes once per second per service, in a background thread so a lost probe (1\,s timeout) never blocks a tick.
\end{itemize}
\end{qa}

\begin{qa}{268}{Trade-off between sampling frequency, detection speed, CPU overhead and noise?}
Faster sampling detects sooner but costs CPU (a \code{tc} fork per read) and measures noisier rates over shorter
intervals. We put fast sampling only where speed matters (link state, as an event) and keep aggregate rates at 1\,s.
\end{qa}

\begin{qa}{269}{Are counters cumulative or instantaneous?}
Kernel interface and HTB counters are cumulative. Backlog is instantaneous. Telemetry converts cumulative counters to
per-tick deltas and rates.
\end{qa}

\begin{qa}{270}{How do you calculate a per-interval byte rate from cumulative counters?}
$r=8\,\max(B_t-B_{t-1},0)/\Delta t/10^6$\,Mb/s, with $\Delta t$ the measured wall-clock interval between reads (not the
nominal 1\,s).
\end{qa}

\begin{qa}{271}{How do you handle counter resets, recreated interfaces, wraparound or a switch restart?}
Negative deltas are clamped to 0 (one zero sample). Sysfs counters are 64-bit, so wraparound is irrelevant. A
missing interface is skipped, its key absent, so the rate reads 0.
\honest{A restart therefore shows as a short zero, not as an explicit ``telemetry gap''.}
\end{qa}

\begin{qa}{272}{What happens when a counter decreases unexpectedly?}
The delta is clamped to zero for that interval and the next interval uses the new baseline. It is not logged as an
anomaly.
\end{qa}

\begin{qa}{273}{How do you calculate utilisation when capacity is known?}
$\max(r_{up},r_{down})/C$: the rural uplink is full-duplex, so it is saturated if either direction is.
\honest{Until the real-OVS validation, utilisation counted only the uplink, which made downlink-heavy contexts (school)
look about 17\% utilised. It was fixed and the whole matrix re-run.}
\end{qa}

\begin{qa}{274}{How do you distinguish offered, forwarded, queued and dropped traffic?}
Offered: access-port rx (uplink) or HTB sent+dropped (downlink). Forwarded: HTB sent. Queued: HTB backlog.
Dropped: HTB drops (queue) and, post-run, matched-packet loss (end to end).
\end{qa}

\begin{qa}{275}{How are per-flow, per-class, per-link and per-slice measurements aggregated?}
Per class = per slice (one HTB class per slice per direction). Per link: interface counters. Per flow: only from
captures, in 5\,s windows. Slice KPIs aggregate flow windows (p95 of per-window p95, mean loss, summed delivered
rate).
\end{qa}

\begin{qa}{276}{Can aggregate utilisation conceal severe congestion on one link?}
Yes, and it did. Uplink-only utilisation hid downlink saturation (Q273). Likewise, a link at 60\% average can still
have one saturated class. That is why decisions use per-class drops and backlog, not just utilisation.
\end{qa}

\begin{qa}{277}{How do you synchronise measurements from different interfaces?}
All readers run in one process within milliseconds on one host clock, and each snapshot carries a timestamp.
Probe results are asynchronous: the snapshot holds the \emph{latest completed} probe (at most about 1\,s old).
\end{qa}

\begin{qa}{278}{How do you handle delayed, missing, duplicated or out-of-order samples?}
Ticks are wall-clock aligned and deltas use measured intervals, so late ticks are handled correctly. A missing probe
becomes \texttt{None}. Duplicates and reordering cannot occur within one process.
\honest{There is no explicit staleness bound on probe results.}
\end{qa}

\begin{qa}{279}{How do you distinguish zero traffic from unavailable telemetry?}
For probes: \texttt{None} means unavailable, distinct from a measured RTT.
\honest{Interface and HTB readers return 0 or empty on error, which is indistinguishable from zero traffic. The fix is
to return \texttt{None} and carry a per-source health flag.}
\end{qa}

\begin{qa}{280}{How do you detect stale telemetry?}
\honest{Not implemented. If the probe thread died, the snapshot would keep reporting its last RTT.}
Fix: timestamp each probe result and treat results older than 2 intervals as missing; add thread liveness checks.
\end{qa}

\begin{qa}{281}{What happens if the monitoring process fails?}
Monitoring runs inside the orchestrator process, so a crash ends the run. The kernel keeps the last QoS configuration
(enforcement survives) and Mininet is torn down by \code{finally}.
\honest{There is no supervisor or restart.}
\end{qa}

\begin{qa}{282}{How do you validate telemetry against independent measurements?}
\begin{itemize}
\item Probe RTT/2 matches the packet-matched one-way delay floor (16\,ms RTT $\leftrightarrow$ 8\,ms one-way).
\item Per-service counter rates match the dataset's delivered throughput (same order per tick).
\item Delay floors match configuration exactly.
\item Unit tests feed synthetic captures with known delays and losses.
\end{itemize}
\end{qa}

\begin{qa}{283}{How do you avoid treating a measurement artefact as a real fault?}
\begin{itemize}
\item Faults are confirmed by physically meaningful signals (carrier loss, sustained trends over a 6-sample window,
congestion evidence).
\item The healer ignores self-inflicted link events within 1\,s of its own actions.
\item Methodological artefacts found during validation were fixed and all runs re-done: busy-wait CPU starvation,
the partial final window, uplink-only utilisation, and false failover from AQM-less queues.
\end{itemize}
\end{qa}

\begin{qa}{284}{Which measurements are unavailable because there is no physical RAN?}
CQI, SINR, RSRP and RSRQ, PRB utilisation, HARQ and BLER, handover events, radio link failures. In the dataset these
columns are scenario-injected placeholders, excluded from ML.
\end{qa}

\section{Feature extraction}

\FigMeasure

\begin{qa}{285}{Why do you need a feature-extraction stage?}
Classifiers and analysis need fixed-size per-flow descriptors over a time window. Raw packets are variable-length,
huge (about 100\,MB of captures per run) and directional. Extraction also produces the QoS ground truth (one-way
delay, loss).
\end{qa}

\begin{qa}{286}{What are the raw inputs and outputs of the feature extractor?}
In: endpoint pcaps (Ethernet or raw IP), UE addresses, spawn metadata (labels), telemetry snapshots.
Out: one record per (application flow, UE, 5\,s window) mapped to the 89-column schema plus measured extras, 31
\code{shape_*} features, and a $3\times32$ packet-sequence tensor (\code{features/extractor.py}).
\end{qa}

\begin{qa}{287}{Which features are computed at Layers 2, 3 and 4?}
\begin{itemize}
\item L2: ether type, VLAN and broadcast flags. These are constant in emulation and kept only for schema compatibility.
\item L3: IP lengths (size statistics), TTL, DSCP.
\item L4: ports, TCP flags (SYN, ACK, FIN, RST), window size, retransmissions, UDP counts.
\end{itemize}
ML uses size, timing, direction and flag counts, but \emph{not} ports, DSCP or TTL (shortcuts).
\end{qa}

\begin{qa}{288}{Which features are flow-level, statistical, temporal or derived?}
\begin{itemize}
\item Flow-level: packet and byte counts per direction, duration.
\item Statistical: size and IAT mean, std, min, max, skewness, kurtosis.
\item Temporal: active and idle time, bursts and burst length, per-direction IAT.
\item Derived: rates, byte share, small-packet share, coefficient of variation.
\item QoS: one-way delay, p95, jitter, loss, queueing delay (labels and analysis only).
\end{itemize}
\end{qa}

\begin{qa}{289}{Why are duration, sizes, rates and inter-arrival statistics useful?}
Services differ in them: ECG and IoT are small, periodic and regular; video is large, bursty and frame-periodic;
VoIP has fixed small packets with talkspurts; bulk TCP is MTU-sized with ACK flows; web is request/response. The top RF
features are forward mean size, forward byte share, backward packet count and burst count.
\end{qa}

\begin{qa}{290}{Which features are available before a flow completes?}
All of them. Records are \emph{windowed} (5\,s), so they exist from the first window onward, and online inference uses
the flow's history so far (at least 5 packets).
\end{qa}

\begin{qa}{291}{Which features require observing the flow over a time window?}
Rates, IAT statistics, burst structure, idle and active time, byte share. The 32-packet sequence needs only the first
32 packets.
\end{qa}

\begin{qa}{292}{How does window length influence accuracy and decision latency?}
Longer windows give more stable statistics but decide later and blur context changes; shorter windows decide earlier
with noisier features.
\honest{Only 5\,s was evaluated. A sweep (1, 2, 5, 10\,s) of F1 against time-to-decision is a clear gap.}
\end{qa}

\begin{qa}{293}{How do you handle short flows with insufficient observations?}
Offline: windows with fewer than 2 packets are dropped. Online: a flow is classified only once it has at least 5
packets; until then it stays in the default (general) class, which is safe but not optimal.
\end{qa}

\begin{qa}{294}{How do you handle missing or invalid feature values?}
Delay is \texttt{NaN} when nothing was delivered in a window (counted by the validator's measurement-coverage score).
Classifiers fill \texttt{NaN} with 0; the CNN applies a sign-log transform and z-scores fitted on training data.
\end{qa}

\begin{qa}{295}{How do you prevent feature definitions diverging between dataset generation and runtime inference?}
One implementation (\code{features/shape.py}) serves both. The model pickle stores its column list and the CNN stores
its normalisation and column order. The ML-feature list lives in \code{registries/features.yaml} and is mirrored in
\code{schema_mapper.py}.
\end{qa}

\begin{qa}{296}{How do you detect unit inconsistencies (bytes vs bits, s vs ms)?}
Units are in every name, and conversions are explicit at one place each (\code{*8/1e6}, \code{*1000}). Unit tests
assert known values (8\,ms one-way delay, 10\% loss, 2\,ms jitter on synthetic captures). The exact propagation floors
are an end-to-end unit check.
\end{qa}

\begin{qa}{297}{How do you handle highly correlated features?}
They are not removed: random forests and CNNs tolerate collinearity, and interpretability claims are limited to
importance rankings.
\honest{The importance of correlated features is split among them, so do not over-interpret individual importances.}
\end{qa}

\begin{qa}{298}{How do you decide whether a feature is informative or a dataset artefact?}
\begin{itemize}
\item Exclude known shortcuts a priori (ports, DSCP, TTL, timestamps, scenario fields, QoS outcomes).
\item Evaluate on held-out run$\times$UE groups.
\item Check that top features are physically meaningful (sizes, direction share, bursts).
\end{itemize}
The strongest available artefact test was the transfer study: synthetic features did not carry over (F1 \NTrZero).
\end{qa}

\begin{qa}{299}{What is the cost of extracting each feature at runtime?}
Shape features are $O(n)$ numpy operations over a flow's packets, well under a millisecond per flow. The cost driver
online is re-parsing the growing UE capture every interval ($O(\text{total packets})$).
\honest{Incremental parsing is needed for long runs.}
\end{qa}

\begin{qa}{300}{Which features risk exposing payload content or sensitive information?}
Features themselves are metadata (sizes, timing), but traffic analysis can still reveal activities.
\honest{The 128\,B capture includes about 60\,B of payload, used only to fingerprint our own generator packets for
delay matching. A real deployment must capture headers only and hash fingerprints.}
\end{qa}

\begin{qa}{301}{How would you run a feature ablation to see which features actually matter?}
Group-wise drop-column importance under group hold-out: remove each family (size, IAT, direction, burst, flags),
retrain, and record $\Delta$F1 with seed variance. Permutation importance gives a quicker check.
\honest{Only RF impurity importances exist today.}
\end{qa}

\begin{qa}{302}{What happens when a feature becomes unavailable after a topology or instrumentation change?}
\honest{Inference re-indexes columns and fills missing ones with 0, which degrades silently.}
Fix: fail loudly when expected features are missing, version the feature set with the model, and retrain.
\end{qa}
